\begin{table}[htbp]
\centering
\footnotesize
\setlength{\tabcolsep}{3.0pt}
\captionstyle{\centering}
\normalcaption
\caption{Inferência alternativa, painel balanceado, contraste entre destinos e ajuste de Holm--Bonferroni}
\label{tab:ap-inferencia-holm}
\begin{tabular}{lcccc}
\toprule
\multicolumn{5}{l}{\textit{Painel A: Inferência alternativa, painel balanceado e contraste Agr. $-$ Past.}} \\
Especificação / Covariância & Est. (Soja) & $p$ (S) & Est. (Milho) & $p$ (M) \\
\midrule
Ref. (Cluster Mun., $G=134$) & 0,0848** (0,0358) & 0,0192 & 0,0519*** (0,0179) & 0,0043 \\
Cluster Duplo (Mun. + Ano) & 0,0848* (0,0481) & 0,0909 & 0,0519* (0,0277) & 0,0741 \\
FWL-HAC (24 escores, lag 2) & 0,0848** (0,0360) & 0,0274 & 0,0519* (0,0252) & 0,0509 \\
Driscoll--Kraay (lag 2) & 0,0848** (0,0382) & 0,0366 & 0,0519* (0,0267) & 0,0647 \\
Painel balanceado ($N=2.256$) & 0,0827** (0,0400) & 0,0413 & 0,0543*** (0,0206) & 0,0098 \\
Empilhada Agr. $-$ Past. ($N=5.534$) & -0,2357** (0,0952) & 0,0145 & -0,4081*** (0,1369) & 0,0034 \\
\midrule
\multicolumn{5}{l}{\textit{Painel B: Correção de multiplicidade de Holm--Bonferroni por família de testes}} \\
Família / Hipótese Testada & Estimativa & EP & $p$ bruto & $p_{\text{Holm}}$ \\
\midrule
\multicolumn{5}{l}{\textbf{F1: H1a e Biomas ($m=6$)}} \\
\quad M1 Soja MT & 0,0848* & (0,0358) & 0,0192 & 0,0767 \\
\quad M2 Milho MT & 0,0519** & (0,0179) & 0,0043 & 0,0248 \\
\quad Soja Amazônia & 0,0866 & (0,0610) & 0,1597 & 0,3194 \\
\quad Soja Cerrado & 0,0752 & (0,0439) & 0,0920 & 0,2759 \\
\quad Milho Amazônia & 0,0185 & (0,0514) & 0,7197 & 0,7197 \\
\quad Milho Cerrado & 0,0572** & (0,0192) & 0,0041 & 0,0248 \\
\multicolumn{5}{l}{\textbf{F2: H1b Instituições ($m=6$)}} \\
\quad Área embargada (Soja) & -0,0076** & (0,0029) & 0,0091 & 0,0454 \\
\quad Área embargada (Milho) & -0,0082** & (0,0030) & 0,0068 & 0,0407 \\
\quad Tripla N embargos (Soja) & -0,0314 & (0,0179) & 0,0810 & 0,3242 \\
\quad Tripla N embargos (Milho) & 0,0068 & (0,0165) & 0,6828 & 1,0000 \\
\quad Autos alinhado (Soja, $N=1.125$) & 0,0034 & (0,0063) & 0,5896 & 1,0000 \\
\quad Tripla autos alinhado (Soja) & -0,0062 & (0,0121) & 0,6084 & 1,0000 \\
\multicolumn{5}{l}{\textbf{F3: Regimes 06--11 e Destinos ($m=6$)}} \\
\quad Painel A Amazônia 2006--11 (Soja) & -0,0643 & (0,0318) & 0,0471 & 0,1247 \\
\quad Painel A Cerrado 2006--11 (Soja) & 0,0523 & (0,0250) & 0,0402 & 0,1247 \\
\quad Painel B Amazônia 2006--11 (Soja) & -0,0704 & (0,0334) & 0,0384 & 0,1247 \\
\quad Painel B Cerrado 2006--11 (Soja) & 0,0569 & (0,0258) & 0,0312 & 0,1247 \\
\quad Empilhado Agr. $-$ Past. (Soja) & -0,2357* & (0,0952) & 0,0145 & 0,0727 \\
\quad Empilhado Agr. $-$ Past. (Milho) & -0,4081** & (0,1369) & 0,0034 & 0,0205 \\
\multicolumn{5}{l}{\textbf{F4: Pecuária / Pastagem ($m=7$)}} \\
\quad Pastagem: Boi $\times H^{\text{past}}$ & -0,0022 & (0,1711) & 0,9897 & 1,0000 \\
\quad Pastagem: Boi $\times$ GAEZ & 0,0040 & (0,1590) & 0,9798 & 1,0000 \\
\quad Pastagem: Razão Soja/Boi $\times$ GAEZ & 0,3941** & (0,1335) & 0,0037 & 0,0161 \\
\quad Pastagem conjunta: Soja $\times$ GAEZ & 0,3963** & (0,1321) & 0,0032 & 0,0161 \\
\quad Pastagem conjunta: Boi $\times$ GAEZ & -0,3413 & (0,2229) & 0,1281 & 0,3843 \\
\quad Pastagem: Soja $\times$ GAEZ & 0,3206** & (0,1016) & 0,0020 & 0,0119 \\
\quad Pastagem: Milho $\times$ GAEZ & 0,4600*** & (0,1397) & 0,0013 & 0,0089 \\
\bottomrule
\end{tabular}
\par\smallskip\begin{minipage}{\textwidth}\footnotesize
\SingleSpacing
\noindent\textbf{Fonte:} elaboração própria. Asteriscos no Painel B indicam significância após ajuste de Holm--Bonferroni (* $p_{\text{Holm}}<0{,}10$; ** $p_{\text{Holm}}<0{,}05$; *** $p_{\text{Holm}}<0{,}01$).\par\medskip
\noindent\textbf{Nota:} No Painel A (asteriscos pelo $p$-valor bruto), soja e milho mantêm significância a 10\% ou menos sob covariâncias conservadoras e no painel contínuo ($N = 2.256$), enquanto o modelo empilhado aponta menor resposta na agricultura que na pastagem. No Painel B, o ajuste de Holm--Bonferroni preserva os efeitos principais (F1, F2 área embargada, F3 destinos e F4 \textit{crop-to-beef}), elevando os contrastes de 2006--2011 para $p_{\text{Holm}} = 0{,}1247$, sem comprovar nulidade.
\end{minipage}
\end{table}
